\documentclass[11pt,a4paper]{article}
\usepackage{fontspec,amsmath,amssymb,amsthm,microtype}
\setmainfont{lmroman10-regular.otf}[BoldFont=lmroman10-bold.otf,ItalicFont=lmroman10-italic.otf,BoldItalicFont=lmroman10-bolditalic.otf]
\usepackage[margin=25mm]{geometry}
\usepackage[unicode,hidelinks]{hyperref}
\hypersetup{pdfauthor={LangevinLean Project}}
\setlength{\emergencystretch}{3em}
\setlength{\parskip}{0.4em}
\setlength{\parindent}{0pt}
\allowdisplaybreaks[2]
\raggedbottom
\title{\texorpdfstring{Construct a stationary Brownian coupling for one true transition}{Construct a stationary Brownian coupling for one true transition}}
\author{LangevinLean Project}
\date{}
\begin{document}\maketitle
Draft mathematical exposition; these notes may contain mistakes.


The analytic one-step premises are realized on the product of an arbitrary finite-cost initial coupling and a canonical Brownian source. The reference is the actual stationary Langevin solution for the smoothed Gibbs potential. Its law, moment bounds and future orthogonality are constructed, while the Brownian endpoint identifies the algorithm's exact independent-Gaussian transition law.




The generic splitting map is \(\mathcal T_h(x,\xi)=T(x-h\nabla f_1(x)+\sqrt{2h}\xi)\). The reference target is the actual Gibbs law of \(F=f_1+e\). Assumptions concern potentials, true energy, bounded gradients and a finite input coupling.



Earlier results used here: \cite[Section~3.45]{LangevinLean2026ULASC}.

\section{\texorpdfstring{Measurable update and exact gradient additivity}{Measurable update and exact gradient additivity}}


The map is the composition of the measurable Gaussian forward update with the Lipschitz map \(T\). Differentiability proves \(\nabla(f_1+e)=\nabla f_1+\nabla e\). Almost everywhere representative identities then prove additivity of the corresponding \(L^2\) drift maps, as an equality of actual \(L^2\) elements.





The smooth forward map is continuous jointly in state and Gaussian input, and the energy map is Lipschitz. Their composition is therefore measurable and defines a transition probability. Differentiability of both potential components gives the sum of their gradients. Composing with square-integrable states respects this equality almost surely, so the reference drift is exactly the score of the sum potential.



Earlier results used here: \cite[Section~3.13.2]{LangevinLean2026ULASC}; \cite[Section~3.13.3]{LangevinLean2026ULASC}; \cite[Section~3.23.2]{LangevinLean2026ULASC}.

\section{\texorpdfstring{Extend any finite-\allowbreak{}cost coupling by canonical Brownian noise}{Extend any finite-cost coupling by canonical Brownian noise}}


Start with any coupling \(c\) between \(\mu\) and the Gibbs target, with finite quadratic cost. Extend it by the canonical independent Brownian source. The initial reference has the exact Gibbs law and finite moments; the constructed Langevin integral solution stays stationary. Its natural-past measurability and Brownian independence prove the future orthogonality. Uniform moments give

\[G=\sqrt{2L^2d/m+2K^2}+K,\qquad\sigma=\sqrt{2d}.\]

Use the Brownian increment over \([0,h]\) as the Gaussian update noise; its independence gives the exact Markov marginal. Apply the derived analytic one-step theorem. The explicit output pair is a finite coupling and proves

\[W_2(\mu\mathcal T_h,\pi)^2\le(1-mh)\operatorname{cost}(c)^2+\mathcal F(m,h,L,G,\sigma,K,\delta).\]

The construction discharges stationarity, Gaussian laws, finite moments and orthogonality; they are not extra convergence hypotheses.




Start from an arbitrary finite-cost coupling \(c\) between the input law and \(\pi\propto e^{-(f_1+e)}\). Extend it by an independent canonical Brownian source, producing square-integrable initial states \(X,Z\) with the required marginals. The actual Lipschitz integral equation constructs \(Y_t\) from \(Z\), and the Gibbs stationarity theorem proves \(\operatorname{Law}(Y_t)=\pi\) throughout the interval.





Under the product source, the first initial marginal has the input law and the second has Gibbs law. Independent Brownian noise evolves the latter by the constructed Langevin equation, whose stationarity theorem preserves the Gibbs law. The normalized endpoint increment is fresh standard Gaussian noise for the former. Their updated pair is thus a finite-cost coupling of the actual transition and the target, to which the derived analytic bound applies.



Earlier results used here: \cite[Section~3.13.2]{LangevinLean2026ULASC}; \cite[Section~3.13.3]{LangevinLean2026ULASC}; \cite[Section~3.17.1]{LangevinLean2026ULASC}; \cite[Section~3.17.12]{LangevinLean2026ULASC}; \cite[Section~3.17.2]{LangevinLean2026ULASC}; \cite[Section~3.17.3]{LangevinLean2026ULASC}; \cite[Section~3.17.4]{LangevinLean2026ULASC}; \cite[Section~3.22.17]{LangevinLean2026ULASC}; \cite[Section~3.22.19]{LangevinLean2026ULASC}; \cite[Section~3.22.2]{LangevinLean2026ULASC}; \cite[Section~3.22.20]{LangevinLean2026ULASC}; \cite[Section~3.22.21]{LangevinLean2026ULASC}; \cite[Section~3.22.22]{LangevinLean2026ULASC}; \cite[Section~3.22.23]{LangevinLean2026ULASC}; \cite[Section~3.22.3]{LangevinLean2026ULASC}; \cite[Section~3.22.7]{LangevinLean2026ULASC}; \cite[Section~3.22.8]{LangevinLean2026ULASC}; \cite[Section~3.25.1]{LangevinLean2026ULASC}; \cite[Section~3.25.4]{LangevinLean2026ULASC}; \cite[Section~3.25.6]{LangevinLean2026ULASC}; \cite[Section~3.26.3]{LangevinLean2026ULASC}; \cite[Section~3.29.7]{LangevinLean2026ULASC}; \cite[Section~3.29.8]{LangevinLean2026ULASC}; \cite[Section~3.32.12]{LangevinLean2026ULASC}; \cite[Section~3.32.13]{LangevinLean2026ULASC}; \cite[Section~3.32.15]{LangevinLean2026ULASC}; \cite[Section~3.32.16]{LangevinLean2026ULASC}; \cite[Section~3.32.17]{LangevinLean2026ULASC}; \cite[Section~3.32.19]{LangevinLean2026ULASC}; \cite[Section~3.33.5]{LangevinLean2026ULASC}; \cite[Section~3.42.6]{LangevinLean2026ULASC}; \cite[Section~3.42.7]{LangevinLean2026ULASC}; \cite[Section~3.8.11]{LangevinLean2026ULASC}; \cite[Section~3.8.16]{LangevinLean2026ULASC}; \cite[Section~3.8.17]{LangevinLean2026ULASC}; \cite[Section~3.8.18]{LangevinLean2026ULASC}; \cite[Section~3.8.19]{LangevinLean2026ULASC}; \cite[Section~3.8.20]{LangevinLean2026ULASC}; \cite[Section~3.8.21]{LangevinLean2026ULASC}.

\begin{thebibliography}{ULA-LSI}
\bibitem[ULA-SC]{LangevinLean2026ULASC}
LangevinLean Project. \emph{Unadjusted Langevin under Strong Convexity}. LL-01, 2026. Draft report, content revision 8d6335dce333. \url{https://odakyildiz.github.io/LangevinLean/docs/pdf/UnadjustedLangevin/FullProof.pdf}.
\end{thebibliography}

\end{document}
